% =====================================================================
%  4. Implementação do abastecimento de água
%  (Atividades 3.2.2.1, 3.2.2.2, 3.2.2.3 e 3.2.2.6 do MOP)
% =====================================================================
\section{Implementação dos sistemas de abastecimento de água}
\label{sec:abastecimento}

A implementação do abastecimento de água percorre seis fases, desde a
preparação institucional até a operação dos sistemas pela entidade gestora
(\autoref{fig:abastecimento}). Em cada comunidade selecionada, o IAT
pactua com o município e com os moradores o modelo de gestão e realiza o
trabalho técnico-social. Em seguida, perfura e outorga o poço, projeta e
licencia o sistema, executa as obras e, por fim, entrega a infraestrutura a
uma entidade capacitada para operá-la. Essa sequência reúne as
\atividade{3.2.2.1} (modelos de gestão), \atividade{3.2.2.2} (trabalho
técnico-social) e \atividade{3.2.2.3} (sistemas coletivos de abastecimento)
do MOP. Ao final da seção, descreve-se a frente complementar do Programa
Sanepar Rural (\atividade{3.2.2.6}), que segue um arranjo próprio.

\begin{figure}[htbp]
  \centering
  \caption{Implementação do abastecimento de água, do começo ao fim}
  \label{fig:abastecimento}
  \fluxograma{fluxograma-abastecimento}
  \fonte{Elaborado pelo IAT (2026) com base em \citeonline{mop2026} e no POA 2026--2027.}
\end{figure}

% ---------------------------------------------------------------------
\subsection{Seleção e adesão dos municípios e comunidades}
\label{subsec:adesao}

A implementação começa dentro do próprio IAT. Entre maio e outubro de 2026,
a equipe técnica do instituto, em conjunto com o IDR-Paraná, fará o
planejamento integrado das ações de saneamento rural a serem executadas a
partir de 2027. Em paralelo, aplicará os critérios de seleção descritos na
\autoref{sec:area} para identificar os municípios elegíveis, cujo resultado
será consolidado em um mapa. Com esse mapa em mãos, o IAT iniciará a
sensibilização das prefeituras por meio de reuniões e ações institucionais,
apresentando os objetivos do Programa e as responsabilidades de cada ente
\cite{mop2026}.

Os municípios interessados formalizarão sua adesão, o que resultará na lista
de municípios participantes, e assinarão com o IAT um Termo de Cooperação.
Esse termo formaliza a parceria: nele ficam registradas as obrigações do Estado, do município e da comunidade durante
toda a implantação, e ele só será encerrado depois que os sistemas forem
entregues e aceitos pela entidade gestora
(\autoref{subsec:entrega}). Assinado o termo, o IAT fará o contato inicial com
as comunidades, que manifestarão seu interesse por meio de contratos de
pré-adesão. As etapas de seleção, adesão e assinatura do termo estão
previstas no POA para o período de maio a dezembro de 2026, sob
responsabilidade da equipe interna do IAT (Anexo~\ref{an:cronograma}). A
\autoref{fig:adesao} mostra como essas etapas se distribuem entre o IAT, o
município e a comunidade.

\begin{figure}[htbp]
  \centering
  \caption{Seleção e adesão dos municípios e comunidades, por responsável}
  \label{fig:adesao}
  \fluxograma{fluxograma-adesao}
  \fonte{Elaborado pelo IAT (2026) com base em \citeonline{mop2026} e no POA 2026--2027.}
\end{figure}

% ---------------------------------------------------------------------
\subsection{Definição e implantação do modelo de gestão}
\label{subsec:modelos-gestao}

O modelo de gestão é o que garante que o sistema continuará funcionando
depois da entrega. A dispersão das comunidades rurais e a limitada capacidade
operacional de muitos municípios tornam necessários arranjos institucionais
que promovam a cooperação entre os diferentes níveis de governo. Os modelos
específicos para o saneamento rural ainda estão em definição. A alternativa
preliminar em análise é um modelo baseado em arranjos intermunicipais,
possivelmente estruturado por meio de consórcios públicos ou mecanismos
equivalentes, que fortaleça a capacidade de gestão dos municípios e garanta
suporte técnico às comunidades. Nesse modelo, Estado, municípios e comunidades
compartilham responsabilidades: o IAT, a SECID e o IDR-Paraná apoiam o
planejamento, a coordenação técnica, a mobilização social e o monitoramento;
os municípios participam da governança e da articulação local; e as
comunidades participam da gestão e da operação cotidiana dos sistemas
\cite{mop2026}. Segundo o POA, a articulação com os entes envolvidos para
definir um modelo de gestão piloto ocorrerá entre abril e dezembro de 2026,
conduzida pelo IAT com o IDR-Paraná e a Sanepar.

Uma vez definido o modelo, sua implantação nas comunidades seguirá seis
etapas, representadas na \autoref{fig:modelo-gestao}. Nos estudos
preliminares, entre 2027 e 2028, serão selecionados os instrumentos de
gestão, definidos os parceiros, critérios e indicadores, escolhidos os meios
de comunicação do modelo e feitos os primeiros contatos com o município e a
comunidade, até a assinatura dos contratos de pré-adesão. Na etapa de
implantação, entre 2028 e 2029, serão elaborados o estatuto, os manuais e as
cartilhas, feito o enquadramento jurídico da proposta, colhida a adesão de
cada família e regularizada a associação que assumirá a gestão. A etapa de
gestão propriamente dita, entre 2030 e 2032, compreende a elaboração dos
manuais de desempenho operacional, os investimentos necessários, a
estruturação do monitoramento e o treinamento das pessoas que vão operar o
sistema. Ao longo de todo esse período, a implantação será acompanhada por
supervisão técnica e avaliada semestralmente, com os dados registrados no
SIGARH e no GeoPR. A atividade se encerra entre 2032 e 2033, com a
consolidação dos resultados, os relatórios finais e a prestação de contas.
Para dar suporte a essa estrutura, o POA prevê ainda a aquisição de
infraestrutura para implantação do modelo de gestão, cujo TR começará a ser
elaborado em novembro de 2027.

\begin{figure}[htbp]
  \centering
  \caption{Etapas de implantação do modelo de gestão e seus produtos}
  \label{fig:modelo-gestao}
  \fluxograma{fluxograma-modelo-gestao}
  \fonte{Elaborado pelo IAT (2026) com base em \citeonline[Quadro 28]{mop2026}.}
\end{figure}

O sucesso desta atividade será medido pelo número de associações de
saneamento rural com modelo de gestão adotado, indicador de resultado do
Componente~3 cuja meta é de 122 associações até 2032
(\autoref{tab:metas-componente3}). O detalhamento das etapas, com os
responsáveis, os produtos, as evidências e a validação de cada uma, está no
\autoref{qua:etapas-3221}, no Apêndice~\ref{ap:quadros}.

\pendencia{No Quadro 28 do MOP (\autoref{qua:etapas-3221}), a etapa
``Supervisão técnica'' repete o produto e as evidências das obras
(``supervisão e gerenciamento de obras'', ``laudos de vistoria e testes de
estanqueidade''). A sigla ``EC'' não está definida. Além disso, o Quadro 31
adota ``tarifa instituída'' como evidência e a SECID como validadora, enquanto
o Quadro 28 usa associações regularizadas e pessoas capacitadas, validadas
pelo IAT. Harmonizar.}

% ---------------------------------------------------------------------
\subsection{Mobilização social e trabalho técnico-social}
\label{subsec:tts}

Com o município e a comunidade formalmente engajados, inicia-se o trabalho
técnico-social, que tem três finalidades: conhecer as famílias e suas
condições de saneamento, definir quem será atendido e validar socialmente as
soluções técnicas propostas. Esse trabalho será executado por empresa
contratada segundo as regras do BIRD. No POA, a mesma contratação reúne a
mobilização social, o levantamento preliminar social, o cadastro das famílias
por busca ativa e os estudos técnicos iniciais do poço (estudo
hidrogeológico, anuência prévia, teste de vazão, análise físico-química e
anteprojeto), descritos na \autoref{subsec:captacao}. O TR dessa
contratação será elaborado entre setembro e novembro de 2026 e a licitação
se estenderá até setembro de 2027, de modo que a execução em campo começará
em outubro de 2027.

Em campo, a empresa começará pela mobilização e sensibilização dos
municípios, das comunidades, das associações e dos proprietários, informando
a população sobre o Programa, estimulando a participação comunitária e
reconhecendo o território onde os sistemas serão implantados. Em seguida,
fará a busca ativa das famílias, inclusive das que vivem em áreas remotas,
acompanhada de ações de capacitação e organização comunitária. Com base
nesse contato, será feito um diagnóstico coletivo das formas atuais de
abastecimento de água e de esgotamento sanitário, que avaliará as condições
existentes, a demanda de água, a infraestrutura necessária e a viabilidade
dos sistemas, complementado pelo levantamento das condições de cada
domicílio.

O diagnóstico permite aplicar os critérios de enquadramento, que combinam
aspectos sociais, territoriais e técnicos, para definir quais famílias serão
atendidas. Cada família enquadrada assinará um termo de adesão, no qual
declara sua participação e seu compromisso com o uso e a manutenção
adequados do sistema. Os resultados serão compilados e apresentados em um
relatório final, que constitui a evidência de conclusão da atividade e dá
baixa na etapa correspondente do Termo de Cooperação. Esse diagnóstico tem
duplo uso: orienta a captação e o projeto do sistema de água e identifica,
domicílio a domicílio, a situação do esgoto, informação de que depende a
frente de esgotamento (\autoref{sec:esgotamento}). A
\autoref{fig:tts} resume a sequência e os produtos.

\begin{figure}[htbp]
  \centering
  \caption{Sequência e produtos do trabalho técnico-social}
  \label{fig:tts}
  \fluxograma{fluxograma-tts}
  \fonte{Elaborado pelo IAT (2026) com base em \citeonline{mop2026} e no POA 2026--2027.}
\end{figure}

\pendencia{O MOP separa o trabalho técnico-social (\atividade{3.2.2.2}) dos
estudos do poço (\atividade{3.2.2.3}), mas o POA os reúne em uma única
contratação. O MOP também coloca a contratação da empresa depois do
diagnóstico, em um item cuja descrição abrange projetos e obras. Harmonizar o
escopo e a sequência.}

% ---------------------------------------------------------------------
\subsection{Captação, projeto e licenciamento}
\label{subsec:captacao}

Definidas as comunidades e as famílias, começam os estudos de engenharia,
que seguem duas linhas encadeadas (\autoref{fig:captacao-projeto}). A primeira
trata da captação subterrânea. Os estudos preliminares avaliarão a
viabilidade técnica, social e ambiental do sistema e verificarão a
titularidade e a situação fundiária das áreas onde ficarão o poço e o
reservatório, cuidado que evita conflitos posteriores sobre o uso da terra.
O levantamento de campo reunirá os dados geológicos e hidrogeológicos
necessários para identificar os aquíferos e as condições ambientais do
local. Com base nele, será definida a locação do ponto de captação e
elaborado o projeto construtivo do poço, que dependerá de anuência prévia
antes da perfuração. Perfurado o poço, serão realizados o teste de vazão e a
análise físico-química da água, que confirmarão se há quantidade e qualidade
suficientes para o abastecimento. Por fim, serão feitos os estudos ambientais
e sociais exigidos para a obtenção da outorga de direito de uso da água
\cite{mop2026}. Para essa linha, o POA prevê a aquisição dos materiais de
perfuração, com TR entre outubro e dezembro de 2026, licitação ao longo de
2027 e início da execução em novembro de 2027.

A segunda linha trata do sistema de tratamento, reservação e distribuição. A
vazão e a qualidade da água obtidas no poço orientam o anteprojeto da
solução de tratamento e armazenamento. Após a ordenação das despesas e a
contratação, serão elaborados o projeto básico e o projeto executivo da casa
química, do reservatório, da parte eletromecânica e da rede de distribuição,
que precisarão ser aprovados antes da obra. Paralelamente, serão realizados
os estudos ambientais e sociais e obtidas as licenças ambientais. O POA prevê
a contratação dos projetos das redes de água e de esgoto com TR entre
novembro de 2026 e janeiro de 2027 e início da execução em dezembro de 2027.
Somente com a outorga e as licenças obtidas as obras poderão começar. O
\autoref{qua:etapas-3223}, no Apêndice~\ref{ap:quadros}, detalha essas etapas
com seus produtos e evidências.

\begin{figure}[htbp]
  \centering
  \caption{Captação subterrânea, projeto e licenciamento do sistema de abastecimento}
  \label{fig:captacao-projeto}
  \fluxograma{fluxograma-captacao-projeto}
  \fonte{Elaborado pelo IAT (2026) com base em \citeonline{mop2026}.}
\end{figure}

\pendencia{O MOP prevê ``contratação integrada'' para o poço e para o sistema
e a define como aquela em que o contratado elabora o projeto executivo e
executa a obra. Na Lei nº~14.133/2021 (art.~6º, XXXII e XXXIII), essa
definição corresponde à contratação \textit{semi-integrada}. Já o POA prevê
outro arranjo: aquisição de materiais pelo IAT e contratação separada dos
projetos. Definir o modelo de execução, indicar quem executará as obras com
os materiais adquiridos e verificar a compatibilidade com o Regulamento de
Aquisições do Banco Mundial.}

% ---------------------------------------------------------------------
\subsection{Execução das obras, supervisão e fiscalização}
\label{subsec:obras}

A obra compreende a construção da casa química, do reservatório, das
instalações eletromecânicas e da rede de distribuição, com as ligações
domiciliares e os hidrômetros. Para essa fase, o POA prevê a aquisição de
materiais para a implantação das redes de abastecimento em dois lotes, de 25
e de 30 unidades. O primeiro terá TR entre dezembro de 2026 e fevereiro de
2027 e licitação até dezembro de 2027; o TR do segundo começará em dezembro
de 2027.

Antes do início das obras, o IAT contratará uma empresa de supervisão e
gerenciamento, que fará o acompanhamento técnico independente da execução.
O mesmo contrato atenderá também às obras de esgotamento sanitário, e seu TR
será elaborado entre dezembro de 2026 e fevereiro de 2027, segundo o POA. A
execução é acompanhada em três níveis, como mostra a \autoref{fig:obras}. O
executor constrói a infraestrutura conforme o projeto. A supervisão
acompanha a obra e registra, em relatórios periódicos, o andamento, os
problemas encontrados e as recomendações. A fiscalização do IAT verifica se
a obra segue o projeto e as normas técnicas e registra essa verificação em
relatórios de fiscalização. Quando a verificação apontar não conformidades,
o executor fará os ajustes necessários antes de nova verificação. Estando
tudo conforme, o IAT emitirá o atesto de recebimento, acompanhado dos laudos
de vistoria, e o executor entregará o projeto \textit{as built}
\cite{mop2026}.

\begin{figure}[htbp]
  \centering
  \caption{Execução, supervisão e fiscalização das obras, por responsável}
  \label{fig:obras}
  \fluxograma{fluxograma-obras}
  \fonte{Elaborado pelo IAT (2026) com base em \citeonline{mop2026}.}
\end{figure}

\pendencia{Os lotes de materiais previstos no POA (25 + 30 = 55 unidades)
não cobrem os 83 sistemas da meta. Confirmar se há lotes posteriores a
2027. A planilha de custos prevê a implantação dos 83 sistemas entre 2028 e
2033; os prazos dos Quadros 29 e 30 do MOP devem ser ajustados a ela. A
quantidade de hidrômetros ainda está como ``XXX''.}

% ---------------------------------------------------------------------
\subsection{Pré-operação, capacitação e entrega à entidade gestora}
\label{subsec:entrega}

Concluída a obra, o sistema passará por um período de pré-operação e
testes, durante o qual os operadores locais serão capacitados. Os
beneficiários também receberão treinamentos sobre operação e manutenção,
comprovados por listas de presença e certificados. A transferência será
formalizada em uma reunião técnica com a comunidade e a entidade gestora, na
qual o IAT entregará o sistema e a entidade confirmará formalmente o seu
recebimento. A partir desse momento, a entidade gestora assume a operação e a
manutenção, enquanto o IAT mantém o suporte técnico pós-implantação e avalia
semestralmente o desempenho dos sistemas por meio de relatórios de
conformidade técnica e ambiental registrados no SIGARH e no GeoPR. A
atividade se encerra com a consolidação dos resultados, o relatório final, a
emissão do Termo de Recebimento Definitivo (TRD) e o encerramento do Termo de
Cooperação com o município \cite{mop2026}. A \autoref{fig:entrega} representa
essa transição.

\begin{figure}[htbp]
  \centering
  \caption{Transferência, operação e encerramento dos sistemas de abastecimento}
  \label{fig:entrega}
  \fluxograma{fluxograma-entrega}
  \fonte{Elaborado pelo IAT (2026) com base em \citeonline{mop2026}.}
\end{figure}

% ---------------------------------------------------------------------
\subsection{Programa Sanepar Rural}
\label{subsec:sanepar-rural}

Nos municípios que mantêm Contrato de Concessão ou de Programa com a Sanepar,
o abastecimento das comunidades rurais será ampliado pelo Programa Sanepar
Rural, com recursos de contrapartida da companhia e execução por convênio
\cite{mop2026}. Estão previstos 222 sistemas coletivos entre 2027 e 2031, com
investimento de R\$ 67,1 milhões \cite{custo2026}. Nesse arranjo, a Sanepar fornece a expertise, os estudos técnicos, o apoio técnico, ambiental e
sociocomunitário, o treinamento, os materiais hidráulicos e os equipamentos,
e cabe ao município executar as obras, com o apoio técnico da companhia. As atividades
seguem o manual do programa \cite{saneparrural}. Começam com o levantamento
técnico e o estudo de engenharia e seguem com a formalização da parceria, por
meio de Termo de Compromisso e Responsabilidade vinculado ao contrato de
concessão. Depois vêm a elaboração dos projetos de engenharia do sistema de
abastecimento, o fornecimento dos materiais e a execução das obras pelo
município. Paralelamente, a comunidade é mobilizada e apoiada na formação de
sua associação de moradores e capacitada em uso racional da água, higiene e
manutenção preventiva. Após a obra, a Sanepar acompanha a qualidade da água e
presta suporte operacional às associações (\autoref{fig:sanepar-rural}). Os
modelos de gestão descritos na \autoref{subsec:modelos-gestao} também se
aplicam a esses sistemas.

\begin{figure}[htbp]
  \centering
  \caption{Programa Sanepar Rural, por responsável}
  \label{fig:sanepar-rural}
  \fluxograma{fluxograma-sanepar-rural}
  \fonte{Elaborado pelo IAT (2026) com base em \citeonline{mop2026}.}
\end{figure}

\pendencia{A meta do Programa Sanepar Rural no Quadro 31 do MOP (59.780
pessoas) é a meta total do indicador ``Pessoas atendidas com água gerida de
forma segura'' do Componente~3, que soma os sistemas do IAT (cerca de 16.200
pessoas) e da Sanepar (cerca de 43.500). A parcela da Sanepar Rural é de
cerca de 43.500 pessoas. O MOP também indica periodicidade ``anual'' no texto e
``contínua'' no Quadro 31, e o manual do programa deve ser citado com título,
versão e data.}

O \autoref{qua:sintese-abastecimento} resume, para cada atividade do
abastecimento, o responsável, o prazo, o produto e a forma de comprovação.

\begin{quadro}[htbp]
\caption{Síntese das atividades de abastecimento de água}
\label{qua:sintese-abastecimento}
\footnotesize
\renewcommand{\arraystretch}{1.25}
\begin{tabularx}{\textwidth}{|P{2.9cm}|P{2.0cm}|Q{1.3cm}|L|L|P{1.4cm}|}
\hline
\cab{Atividade} & \cab{Responsável} & \cab{Prazo} & \cab{Produto} &
\cab{Evidência} & \cab{Validação} \\ \hline
3.2.2.1 Modelos de gestão & IAT, IDR-Paraná e Sanepar & 2027 a 2033 &
Gestão dos sistemas implantada & Tarifa instituída & SECID \\ \hline
3.2.2.2 Trabalho técnico-social & IAT e IDR-Paraná & 2027 a 2030 &
Diagnóstico e cadastro domiciliar & Relatório final e baixa no Termo de
Cooperação & IAT e IDR-Paraná \\ \hline
3.2.2.3 Sistemas coletivos de abastecimento & IAT e IDR-Paraná & 2028 a 2033 &
83 sistemas implantados & Termo de Recebimento Definitivo & IAT \\ \hline
3.2.2.6 Programa Sanepar Rural & Sanepar e IAT & 2027 a 2031 &
222 sistemas implantados (pessoas atendidas) & Relatório técnico aprovado &
Sanepar \\ \hline
\end{tabularx}
\fonte{Adaptado de \citeonline[Quadros 30 e 31]{mop2026}; prazos conforme \citeonline{custo2026}.}
\end{quadro}
